\documentclass{article}
\usepackage{amsmath,amssymb,amsthm}
\usepackage{bm}
\usepackage{hyperref}
\usepackage{enumitem}
\newtheorem{theorem}{Theorem}
\newtheorem{lemma}{Lemma}
\newcommand{\norm}[1]{\left\lVert #1 \right\rVert}
\newcommand{\inner}[2]{\left\langle #1,#2 \right\rangle}
\begin{document}

\section*{Algorithm Name}
\textbf{Gradient Descent–Ascent (GDA) for a Min–Max Problem}  

We consider the saddle‑point problem  

\[
\min_{y\in\mathbb R^{d_y}}\;\max_{x\in\mathbb R^{d_x}} \;L(x,y),
\]

where the objective function \(L:\mathbb R^{d_x}\times\mathbb R^{d_y}\to\mathbb R\) is smooth, possibly non‑convex in \(x\) and concave in \(y\).

\section*{Mathematical Formulation}
Let the iterates be \(x_t\in\mathbb R^{d_x}\) and \(y_t\in\mathbb R^{d_y}\).  
For a stepsize \(\eta>0\) the **GDA update** is  

\[
\boxed{
\begin{aligned}
x_{t+1} &= x_t + \eta \,\nabla_x L(x_t,y_t) , \qquad\text{(ascent in }x\text{)}\\[4pt]
y_{t+1} &= y_t - \eta \,\nabla_y L(x_t,y_t) , \qquad\text{(descent in }y\text{)} .
\end{aligned}}
\tag{1}
\]

All quantities are evaluated at the same time index \(t\); the same scalar stepsize \(\eta\) is used for both players.

\section*{Pseudocode}
\begin{enumerate}[label=\textbf{Step \arabic*:}, leftmargin=*, itemsep=2pt]
\item \textbf{Input:} initial point \((x_0,y_0)\), stepsize \(\eta>0\), maximum number of iterations \(T_{\max}\), tolerance \(\varepsilon\).
\item \textbf{For} \(t=0,1,\dots ,T_{\max}-1\) \textbf{do}
\begin{enumerate}[label=\arabic*.]
\item Compute stochastic (or exact) partial gradients
\[
g^x_t \gets \nabla_x L(x_t,y_t),\qquad
g^y_t \gets \nabla_y L(x_t,y_t).
\]
\item Update the primal (max) variable
\[
x_{t+1} \gets x_t + \eta\, g^x_t .
\]
\item Update the dual (min) variable
\[
y_{t+1} \gets y_t - \eta\, g^y_t .
\]
\item If \(\norm{g^x_t}^2+\norm{g^y_t}^2\le\varepsilon\) then \textbf{break}.
\end{enumerate}
\item \textbf{Output:} the last iterate \((x_T,y_T)\) and optionally the averaged iterate 
\(\bar x_T=\tfrac1T\sum_{t=0}^{T-1}x_t,\;\bar y_T=\tfrac1T\sum_{t=0}^{T-1}y_t\).
\end{enumerate}

\section*{Dry Run Example}
We illustrate the update rule on the simple scalar function  

\[
f(w)= (w-3)^2 ,\qquad w_0=0,\qquad \eta =0.1 .
\]

Because there is only one variable, the algorithm reduces to ordinary gradient descent:
\(w_{t+1}=w_t-\eta \nabla f(w_t)\).

\[
\nabla f(w)=2(w-3).
\]

\begin{center}
\begin{tabular}{c|c|c|c}
\textbf{Iter.} & $w_t$ & $\nabla f(w_t)$ & $w_{t+1}=w_t-\eta\nabla f(w_t)$ \\ \hline
0 & $0$ & $2(0-3)=-6$ & $0-0.1(-6)= -0.6$ \\[4pt]
1 & $-0.6$ & $2(-0.6-3)=-7.2$ & $-0.6-0.1(-7.2)= 0.12$ \\[4pt]
2 & $0.12$ & $2(0.12-3)=-5.76$ & $0.12-0.1(-5.76)=0.696$ \\[4pt]
\end{tabular}
\end{center}

The corresponding objective values are  

\[
f(-0.6)= ( -0.6-3)^2 = 12.96,\qquad 
f(0.12)= (0.12-3)^2 = 8.2944,\qquad 
f(0.696)= (0.696-3)^2 = 5.299\; .
\]

The objective strictly decreases, illustrating the behaviour of the descent step (the ascent step would be analogous with a sign flip).

\newpage
\section*{Assumptions}
\begin{enumerate}[label=\textbf{A\arabic*:}, leftmargin=*]
\item \textbf{Smoothness.}  
\(L\) is \(L\)-smooth in the joint variable \(z=(x,y)\); i.e. for all \(z,z'\),

\[
\| \nabla L(z)-\nabla L(z')\| \le L\|z-z'\|.
\tag{A1}
\]

Equivalently,
\[
L(z')\le L(z)+\inner{\nabla L(z)}{z'-z}+\frac{L}{2}\|z'-z\|^{2}.
\tag{A1'}
\]

\item \textbf{Bounded Level Set.}  
The initial point \((x_0,y_0)\) lies in a level set of finite diameter:
there exists \(D>0\) such that \(\|z-z^*\|\le D\) for all iterates, where
\(z^*=(x^*,y^*)\) denotes any saddle point (existence assumed).

\item \textbf{Existence of a Saddle Point.}  
There exists \((x^*,y^*)\) satisfying  

\[
L(x^*,y)\le L(x^*,y^*)\le L(x,y^*)\quad\forall x,y .
\tag{A3}
\]

\item \textbf{Stepsize.}  
The stepsize satisfies  

\[
0<\eta\le \frac{1}{L}.
\tag{A4}
\]

\end{enumerate}

\section*{Main Convergence Theorem}
\begin{theorem}[Stationarity of GDA]\label{thm:main}
Assume (A1)–(A4).  Let the iterates \(\{(x_t,y_t)\}_{t\ge0}\) be generated by \eqref{1}.  
Define the \emph{game gradient}
\[
G(z_t):=\begin{pmatrix}
\nabla_x L(x_t,y_t)\\[2pt]
-\nabla_y L(x_t,y_t)
\end{pmatrix},\qquad z_t:=(x_t,y_t).
\]

Then for any integer \(T\ge1\),

\[
\frac{1}{T}\sum_{t=0}^{T-1}\|G(z_t)\|^{2}
\;\le\;
\frac{2\bigl(L(x_0,y_0)-L^*\bigr)}{\eta\,T},
\tag{2}
\]

where \(L^*:=\inf_{x}\sup_{y}L(x,y)=L(x^*,y^*)\).  Consequently,
\[
\min_{0\le t<T}\|G(z_t)\|^{2}=O\!\left(\frac{1}{T}\right).
\]
\end{theorem}

\section*{Theoretical Properties}
\begin{enumerate}[label=\textbf{P\arabic*:}, leftmargin=*]
\item \textbf{Stability.}  
Because the update is a standard gradient step on the \emph{monotone operator} \(G\) with Lipschitz constant \(L\), the condition \(\eta\le 1/L\) guarantees that the potential \(L(x_t,y_t)\) does not increase unboundedly; the iterates remain inside the bounded level set (A2).

\item \textbf{Hyper‑parameter Sensitivity.}  
The bound \eqref{2} scales linearly with \(1/\eta\).  A larger stepsize accelerates the decrease of the right‑hand side but must respect \(\eta\le 1/L\).  If \(\eta\) is chosen too large, the term \((1-\tfrac{L\eta}{2})\) in the descent inequality becomes negative and the guarantee fails.

\item \textbf{Comparison to Baselines.}  
For pure gradient descent on a non‑convex function the same \(O(1/T)\) rate for the average gradient norm is classical.  The theorem shows that, despite the opposing ascent direction in \(x\), GDA enjoys the identical rate under the smoothness assumption, matching the best known rates for first‑order methods in non‑convex optimisation.

\item \textbf{Special Cases.}
\begin{itemize}
\item If \(L\) is convex in \(x\) and concave in \(y\) (the classical saddle‑point setting), the same bound holds and, together with a primal–dual gap argument, yields an \(O(1/T)\) convergence of the gap.
\item If the stepsize is chosen as \(\eta_t = \frac{c}{\sqrt{t+1}}\) (diminishing), the bound improves to \(O(\frac{\log T}{\sqrt{T}})\) for the averaged gradient norm (standard extension).
\end{itemize}
\end{enumerate}

\section*{Preliminaries (Definitions and Lemmas)}
\begin{lemma}[Smoothness inequality]\label{lem:smooth}
If \(L\) satisfies (A1), then for any \(z,z'\in\mathbb R^{d_x+d_y}\),

\[
L(z')\le L(z)+\inner{\nabla L(z)}{z'-z}+\frac{L}{2}\|z'-z\|^{2}.
\tag{3}
\]
\end{lemma}
\begin{proof}
Immediate from the definition of an \(L\)-Lipschitz gradient (integrate the gradient bound). \hfill\(\square\)
\end{proof}

\begin{lemma}[Descent Lemma for the Game Gradient]\label{lem:descent}
Let \(z_{t+1}=z_t-\eta G(z_t)\) with \(\eta\le 1/L\).  Then

\[
L(z_{t+1})\le L(z_t)-\frac{\eta}{2}\|G(z_t)\|^{2}.
\tag{4}
\]
\end{lemma}
\begin{proof}
Apply Lemma~\ref{lem:smooth} with \(z'=z_{t+1}=z_t-\eta G(z_t)\):
\[
\begin{aligned}
L(z_{t+1}) 
&\le L(z_t)+\inner{\nabla L(z_t)}{-\eta G(z_t)}+\frac{L}{2}\|\eta G(z_t)\|^{2}\\[4pt]
&= L(z_t)-\eta\inner{\nabla L(z_t)}{G(z_t)}+\frac{L\eta^{2}}{2}\|G(z_t)\|^{2}.
\end{aligned}
\tag{5}
\]
Recall the definition of \(G\):
\[
G(z_t)=\begin{pmatrix}\nabla_x L(z_t)\\ -\nabla_y L(z_t)\end{pmatrix},
\qquad
\nabla L(z_t)=\begin{pmatrix}\nabla_x L(z_t)\\ \nabla_y L(z_t)\end{pmatrix}.
\]
Hence \(\inner{\nabla L(z_t)}{G(z_t)}=\|\nabla_x L(z_t)\|^{2}+\|\nabla_y L(z_t)\|^{2}
=\|G(z_t)\|^{2}\).  Substituting into \eqref{5} yields  

\[
L(z_{t+1})\le L(z_t)-\eta\|G(z_t)\|^{2}+\frac{L\eta^{2}}{2}\|G(z_t)\|^{2}
= L(z_t)-\Bigl(\eta-\frac{L\eta^{2}}{2}\Bigr)\|G(z_t)\|^{2}.
\tag{6}
\]

Because \(\eta\le 1/L\) we have \(\eta-\tfrac{L\eta^{2}}{2}\ge \tfrac{\eta}{2}\).  Therefore  

\[
L(z_{t+1})\le L(z_t)-\frac{\eta}{2}\|G(z_t)\|^{2},
\]
which is exactly \eqref{4}. \hfill\(\square\)
\end{proof}

\section*{Main Proof (Step‑by‑Step Derivation)}
We now prove Theorem~\ref{thm:main} using Lemmas~\ref{lem:smooth} and \ref{lem:descent}.

\begin{enumerate}[label=\textbf{[Step \arabic*]}, leftmargin=*]
\item \textbf{Apply the descent inequality} \eqref{4} to each iteration \(t=0,\dots ,T-1\):
\[
L(z_{t+1})\le L(z_t)-\frac{\eta}{2}\|G(z_t)\|^{2}.
\tag{7}
\]

\item \textbf{Sum telescopically} over \(t=0\) to \(T-1\):
\[
\sum_{t=0}^{T-1}\bigl(L(z_t)-L(z_{t+1})\bigr)
\ge \frac{\eta}{2}\sum_{t=0}^{T-1}\|G(z_t)\|^{2}.
\tag{8}
\]

\item \textbf{Recognise the telescoping sum} on the left‑hand side:
\[
L(z_0)-L(z_T)\ge \frac{\eta}{2}\sum_{t=0}^{T-1}\|G(z_t)\|^{2}.
\tag{9}
\]

\item \textbf{Lower bound the terminal value} \(L(z_T)\) by the saddle‑point value \(L^*\).  
Since \(L^*\) is the infimum over \(x\) of the supremum over \(y\), we have for any iterate
\(L(z_T)\ge L^*\).  Hence

\[
L(z_0)-L^*\ge \frac{\eta}{2}\sum_{t=0}^{T-1}\|G(z_t)\|^{2}.
\tag{10}
\]

\item \textbf{Divide by \(T\)} and rearrange to obtain the average bound:

\[
\frac{1}{T}\sum_{t=0}^{T-1}\|G(z_t)\|^{2}
\le \frac{2\bigl(L(z_0)-L^*\bigr)}{\eta\,T}.
\tag{11}
\]

\item \textbf{Conclusion.}  Inequality \eqref{11} is exactly the statement \eqref{2} of Theorem~\ref{thm:main}.  Moreover,
\(\min_{0\le t<T}\|G(z_t)\|^{2}\le \tfrac{1}{T}\sum_{t=0}^{T-1}\|G(z_t)\|^{2}=O(1/T)\), establishing the claimed convergence rate. \hfill\(\square\)
\end{enumerate}

\section*{Rate Derivation (With Constants)}
From \eqref{11} we can read off the explicit constant:

\[
C:=\frac{2\bigl(L(x_0,y_0)-L^*\bigr)}{\eta}.
\]

Thus for any \(T\ge1\),

\[
\min_{0\le t<T}\|G(z_t)\|^{2}\le \frac{C}{T}.
\]

If a diminishing stepsize \(\eta_t=\frac{c}{\sqrt{t+1}}\) is employed, a standard summation argument gives  

\[
\frac{1}{T}\sum_{t=0}^{T-1}\|G(z_t)\|^{2}
\le \frac{2\bigl(L(x_0,y_0)-L^*\bigr)}{c}\,\frac{\log T}{\sqrt{T}} .
\]

\section*{Special Cases}
\begin{enumerate}[label=\arabic*.]
\item \textbf{Convex–Concave Setting.}  
If \(L\) is convex in \(x\) and concave in \(y\), the saddle point is unique and the primal–dual gap satisfies  

\[
\max_{y}\!L(\bar x_T,y)-\min_{x}\!L(x,\bar y_T)\le \frac{C}{T}.
\]

\item \textbf{Strongly Concave in \(y\).}  
Assume \(-L(\cdot,y)\) is \(\mu\)-strongly convex in \(y\).  Then the dual iterates converge linearly to \(y^*\) while the primal iterates retain the \(O(1/T)\) bound for the gradient norm.

\item \textbf{Stochastic Gradients.}  
If only unbiased estimates \(\mathbb{E}[g^x_t|\mathcal F_t]=\nabla_x L\) and \(\mathbb{E}[g^y_t|\mathcal F_t]=\nabla_y L\) with bounded variance \(\sigma^2\) are available, the same derivation yields  

\[
\mathbb{E}\!\left[\frac{1}{T}\sum_{t=0}^{T-1}\|G(z_t)\|^{2}\right]
\le \frac{2(L(z_0)-L^*)}{\eta T}+ \eta \sigma^2 .
\]

Choosing \(\eta = \Theta(T^{-1/2})\) balances the two terms and leads to the optimal stochastic rate \(O(T^{-1/2})\).
\end{enumerate}

\section*{Discussion}
The proof shows that GDA is nothing more than a gradient step on the \emph{game gradient} \(G\).  Under the sole smoothness assumption the standard descent lemma yields an \(O(1/T)\) decay of the squared norm of \(G\).  This matches the optimal rate for first‑order methods on non‑convex optimisation and demonstrates that the presence of an ascent direction does not deteriorate the convergence speed, provided the stepsize respects the Lipschitz constant.  Extensions to stochastic gradients, diminishing stepsizes, and additional curvature (strong concavity in \(y\)) follow by elementary modifications of the same argument.

\end{document}